% GENERATED FILE. Edit canonical lessons, then run npm run generate:books.
% canonical-source-hash: fnv1a64:57ef9c83c0378ff2
% canonical-lessons: ML-C153-pravu, ML-C153-pumparra, ML-C153-urumpu, ML-C153-kotuku, ML-C153-tol

\chapter{Dove, Butterfly, Ant, Mosquito, Shoulder}
\label{ch:ml-dove-butterfly-ant-mosquito-shoulder}
\hlchaptermodality{153}

\begin{chapteropening}
\textbf{By the end of this chapter:} \emph{I can say a dove, a butterfly, an ant, a mosquito and a shoulder.}

Complete the last lesson of chapter 153: i can say a dove, a butterfly, an ant, a mosquito and a shoulder.
\end{chapteropening}

\section[prāvŭ]{\ml{പ്രാവ്} (prāvŭ) — a dove, a pigeon}

\label{lesson:ML-C153-pravu}



\begin{quote}
Before the new one: say the Malayalam for a peacock, then the Malayalam for a parrot.

Which \textquotedblleft{}day\textquotedblright{} ending do \emph{sāyāhnaṁ} and \emph{madhyāhnaṁ} share? (\emph{-ahna}.)
\end{quote}

\subsection*{You'll want to know: \ml{പ്രാവ്}}
\textbf{\ml{പ്രാവ്}} — \emph{prāvŭ} — \textquotedblleft{}a dove, a pigeon\textquotedblright{}.

\begin{grammarlens}[title={What you've built}]
One more word for this chapter.
\end{grammarlens}

\subsection*{Your turn}
\begin{itemize}
\raggedright
  \item \emph{Say it:} \emph{prāvŭ}
  \item \emph{Say it:} \emph{prāvŭ}, once more
  \item \emph{Say it:} say \emph{prāvŭ}
  \item \emph{Recall:} say the Malayalam for a peacock, then the Malayalam for a parrot, then say \emph{prāvŭ} again
\end{itemize}

\begin{tcolorbox}[breakable,colback=teal!4,colframe=teal!35!black,title={Before you move on}]
What does \ml{പ്രാവ്} mean? (\textquotedblleft{}A dove, a pigeon\textquotedblright{}.) Say it once more.
\end{tcolorbox}
\section[pūmpāṟṟa]{\ml{പൂമ്പാറ്റ} (pūmpāṟṟa) — a butterfly}

\label{lesson:ML-C153-pumparra}



\begin{quote}
Before the new one: say the Malayalam for a parrot, then the Malayalam for a dove.
\end{quote}

\subsection*{You'll want to know: \ml{പൂമ്പാറ്റ}}
\textbf{\ml{പൂമ്പാറ്റ}} — \emph{pūmpāṟṟa} — \textquotedblleft{}a butterfly\textquotedblright{}.

\begin{grammarlens}[title={What you've built}]
One more word for this chapter.
\end{grammarlens}

\subsection*{Your turn}
\begin{itemize}
\raggedright
  \item \emph{Say it:} \emph{pūmpāṟṟa}
  \item \emph{Say it:} \emph{pūmpāṟṟa}, once more
  \item \emph{Say it:} say \emph{pūmpāṟṟa}
  \item \emph{Recall:} say the Malayalam for a parrot, then the Malayalam for a dove, then say \emph{pūmpāṟṟa} again
\end{itemize}

\begin{tcolorbox}[breakable,colback=teal!4,colframe=teal!35!black,title={Before you move on}]
What does \ml{പൂമ്പാറ്റ} mean? (\textquotedblleft{}A butterfly\textquotedblright{}.) Say it once more.
\end{tcolorbox}
\section[uṟumpŭ]{\ml{ഉറുമ്പ്} (uṟumpŭ) — an ant}

\label{lesson:ML-C153-urumpu}



\begin{quote}
Before the new one: say the Malayalam for a dove, then the Malayalam for a butterfly.

Say the formal \textquotedblleft{}good afternoon,\textquotedblright{} the one three languages build the same way. (\emph{Śubha madhyāhnaṁ}.)
\end{quote}

\subsection*{You'll want to know: \ml{ഉറുമ്പ്}}
\textbf{\ml{ഉറുമ്പ്}} — \emph{uṟumpŭ} — \textquotedblleft{}an ant\textquotedblright{}.

\begin{grammarlens}[title={What you've built}]
One more word for this chapter.
\end{grammarlens}

\subsection*{Your turn}
\begin{itemize}
\raggedright
  \item \emph{Say it:} \emph{uṟumpŭ}
  \item \emph{Say it:} \emph{uṟumpŭ}, once more
  \item \emph{Say it:} say \emph{uṟumpŭ}
  \item \emph{Recall:} say the Malayalam for a dove, then the Malayalam for a butterfly, then say \emph{uṟumpŭ} again
\end{itemize}

\begin{tcolorbox}[breakable,colback=teal!4,colframe=teal!35!black,title={Before you move on}]
What does \ml{ഉറുമ്പ്} mean? (\textquotedblleft{}An ant\textquotedblright{}.) Say it once more.
\end{tcolorbox}
\section[kotukŭ]{\ml{കൊതുക്} (kotukŭ) — a mosquito}

\label{lesson:ML-C153-kotuku}



\begin{quote}
Before the new one: say the Malayalam for a butterfly, then the Malayalam for an ant.
\end{quote}

\subsection*{You'll want to know: \ml{കൊതുക്}}
\textbf{\ml{കൊതുക്}} — \emph{kotukŭ} — \textquotedblleft{}a mosquito\textquotedblright{}.

\begin{grammarlens}[title={What you've built}]
One more word for this chapter.
\end{grammarlens}

\subsection*{Your turn}
\begin{itemize}
\raggedright
  \item \emph{Say it:} \emph{kotukŭ}
  \item \emph{Say it:} \emph{kotukŭ}, once more
  \item \emph{Say it:} say \emph{kotukŭ}
  \item \emph{Recall:} say the Malayalam for a butterfly, then the Malayalam for an ant, then say \emph{kotukŭ} again
\end{itemize}

\begin{tcolorbox}[breakable,colback=teal!4,colframe=teal!35!black,title={Before you move on}]
What does \ml{കൊതുക്} mean? (\textquotedblleft{}A mosquito\textquotedblright{}.) Say it once more.
\end{tcolorbox}
\section[tōḷ]{\ml{തോൾ} (tōḷ) — a shoulder}

\label{lesson:ML-C153-tol}



\begin{quote}
Before the new one: say the Malayalam for an ant, then the Malayalam for a mosquito.

Say \emph{śubha madhyāhnaṁ} once more, the formal afternoon greeting the family shares.
\end{quote}

\subsection*{You'll want to know: \ml{തോൾ}}
\textbf{\ml{തോൾ}} — \emph{tōḷ} — \textquotedblleft{}a shoulder\textquotedblright{}.

\begin{grammarlens}[title={What you've built}]
One more word for this chapter.
\end{grammarlens}

\subsection*{Your turn}
\begin{itemize}
\raggedright
  \item \emph{Say it:} \emph{tōḷ}
  \item \emph{Say it:} \emph{tōḷ}, once more
  \item \emph{Say it:} say \emph{tōḷ}
  \item \emph{Recall:} say the Malayalam for an ant, then the Malayalam for a mosquito, then say \emph{tōḷ} again
\end{itemize}

\begin{tcolorbox}[breakable,colback=teal!4,colframe=teal!35!black,title={Before you move on}]
What does \ml{തോൾ} mean? (\textquotedblleft{}A shoulder\textquotedblright{}.) Say it once more.
\end{tcolorbox}
